\slajdid{02-s018}
\begin{frame}[label=02-s018]
\gusttekst
\begin{columns}[c,onlytextwidth]
\begin{column}{.62\textwidth}
Oba načina prikazivanja funkcija možemo sada zapisati i na jednostavniji način korišćenjem definisanih indeksa proizvoda odnosno zbirova. Jasno je sada da ti indeksi predstavljaju vrednosti binarnih brojeva, kada stanja ulaznih promenljivih tretiramo kao binarne cifre, u odgovarajućoj vrsti. Znači
\[
F=\sum\nolimits_{C,B,A}(3,5,6)
\]
\[
F=\prod\nolimits_{C,B,A}(0,1,2,4,7)
\]
„cifra C je najveće težine, pa cifra B pa cifra A“
\end{column}
\begin{column}{.35\textwidth}
\begingroup
\setlength{\tabcolsep}{3pt}
\renewcommand{\arraystretch}{1.2}
\par\noindent
\begin{tabularx}{\linewidth}{l*{4}{>{\centering\arraybackslash}X}}
\toprule indeks&C&B&A&F\\\midrule
0&0&0&0&0\\
1&0&0&1&0\\
2&0&1&0&0\\
3&0&1&1&1\\
4&1&0&0&0\\
5&1&0&1&1\\
6&1&1&0&1\\
7&1&1&1&0\\\bottomrule
\end{tabularx}\par
\endgroup

\end{column}
\end{columns}
\centering
Vrednost binarnog broja \(b_{n-1}b_{n-2}b_{n-3}\ldots b_1b_0\)
\[
\alert{B=b_{n-1}2^{n-1}+b_{n-2}2^{n-2}+b_{n-3}2^{n-3}+\ldots+b_12^1+b_02^0}
\]
\end{frame}
